% Source: https://artofproblemsolving.com/community/c3h2026921_moldy_cheese_mock_mathcounts_problems_1115
% Stable identity: aops-topic-2026921-q15
The six points are the vertices of a regular hexagon inscribed in the unit circle. Their midpoints with the centre $O$ lie at distance $\dfrac{1}{2}$ from $O$, equally spaced at $60^\circ$, so they form a regular hexagon with circumradius $\dfrac{1}{2}$.
\begin{asy}
usepackage("amsmath");
size(4.5cm);
defaultpen(fontsize(10pt));
pair O=(0,0);
draw(unitcircle,gray);
pair[] P, M;
for(int i=0;i<6;++i){ P[i]=dir(60*i); M[i]=P[i]/2; draw(O--P[i],gray+dashed); dot(P[i]); }
draw(M[0]--M[1]--M[2]--M[3]--M[4]--M[5]--cycle,blue+linewidth(0.9));
fill(O--M[0]--M[1]--cycle,paleblue);
draw(M[0]--M[1]--M[2]--M[3]--M[4]--M[5]--cycle,blue+linewidth(0.9));
dot(O); label("$O$",O,SW);
label("$1$",P[2]*0.75,W);
label("$\dfrac{1}{2}$",(M[0]+M[1])/2,NE,blue);
label("\mbox{}",(min(currentpicture,user=true)+max(currentpicture,user=true))/2,black); // keeps the text after the figure black when LaTeX typesets the labels inline
\end{asy}

That hexagon splits into $6$ equilateral triangles with side $\dfrac{1}{2}$ (each has two sides of length $\dfrac{1}{2}$ through $O$ with a $60^\circ$ angle between them). Each has area $\dfrac{\sqrt3}{4}\cdot\dfrac{1}{4}=\dfrac{\sqrt3}{16}$, so the hexagon's area is
\[6\cdot\dfrac{\sqrt3}{16}=\boxed{\dfrac{3\sqrt3}{8}}.\]
